\section{诉讼时效}
\begin{enumerate}
  \item 效力法定。催促行权期间。
        \begin{enumerate}
          \item 预先放弃无效。\al{合同约定}。
          \item 事后放弃（\ydy 不行权）。部分同意（单方允诺之债）则部分重新起算。\al{甲不知诉讼时效已过，向乙还钱}。
          \item 当事人不能约定延长 \hz 缩短。
        \end{enumerate}
  \item 仅适用于请求权（财产性请求权居多）。可以并用。
        \begin{enumerate}
          \item 债权请求权。要钱不要物。1 个原则（3 年），5 个例外。存、债、资、国、费。
                \begin{enumerate}
                  \item 存款本息请求权、债券本息请求权。
                  \item 基于投资关系而产生的缴付出资请求权。人身性。
                  \item 国家财产受到侵犯。防止国有资产流失。
                  \item 3 费。抚养费、扶养费、赡养费请求权。生存权。人身性。这个考的相对多点。
                \end{enumerate}
          \item 物权请求权。要物不要钱。
                \begin{enumerate}
                  \item 1 个原则（不受）。请求停止侵害、排除妨害、消除危险、不动产物权、登记的动产返还。常考。
                  \item 2 个例外。普通动产返还（3 年）、抵押权人行权。
                \end{enumerate}
        \end{enumerate}
  \item 届满后。实体、程序权利均不消灭。
        \begin{enumerate}
          \item 实体权利不消灭，转为自然之债。\al{诉讼时效届满后所受领的给付构成不当得利 \cuo}。
          \item 诉讼权利不消灭，仍可诉且需受理。结合民诉，应当受理，但驳回起诉。
          \item 抗辩权发生说。义务人取得拒绝履行的抗辩权（诉讼时效抗辩权）。
        \end{enumerate}
  \item 种类。
        \begin{enumerate}
          \item 普通/一般诉讼时效。3 年。可变期间（可中止、中断）。
          \item 特殊诉讼时效。涉外、技术进口 \yjt 4 年。人寿保险 \yjt 5 年。
          \item 最长保护期限。20 年。自权利受到损害之日起（客观起算）。可依申请延长。\al{医院抱错小孩}。
        \end{enumerate}
  \item 起算。主观知道权利受到损害以及义务人之日起。例外。
        \begin{enumerate}
          \item 合同之债。\al{典型如借款}。
                \begin{enumerate}
                  \item 未约定履行期限 \yjt 随时行权。
                  \item 约定 \yjt（最后一期）履行期限届满之日。
                  \item \zd{连带保证互不干涉原则}！！连带保证人责任太重了。
                \end{enumerate}
        \end{enumerate}
  \item 中止（\bs{194}）。\zx{最后 6 个月内，因客观障碍，导致不能行使请求权的，诉讼时效中止}。考的少。
        \begin{enumerate}
          \item 不可抗力。
          \item 无限民无法代。
          \item 未确定继承人或者遗产管理人。
          \item 权利人被义务人或者其他人控制。
          \item \bs{婚姻关系存续期间，当事人不起诉离婚而单独提起损害赔偿请求的，不予受理}。
        \end{enumerate}
  \item \zd{中断。因权利人主张权利而重新起算}。
  \item 绝对地说二审期间可以/不可以提出诉讼时效的抗辩，都是不正确的。
  \item vs 除斥期间。除斥期间对应形成权（单方意思表示就可使民事法律关系发生变动）。
        \begin{enumerate}
          \item 适用对象不同。请求权/形成权。
          \item 构成要件不同。法定期间经过 + 权利继续不行使的事实状态/法定期间经过。
          \item 法律效力不同。不行使，权利本身不消灭，而只是消灭附着于其上的胜诉权/权利本身消灭。
          \item 期间起算点不同。主观起算/权利成立之时起算。
          \item 期间是否可变不同。可变期间/不变期间。
          \item 法院是否可依职权主动适用不同。经权利人主张/应依职权主动调查适用。
        \end{enumerate}
\end{enumerate}
